\documentclass{amsart}
% For dealing with references we use the comment environment
\usepackage{verbatim}
\newenvironment{reference}{\comment}{\endcomment}
%\newenvironment{reference}{}{}
\newenvironment{slogan}{\comment}{\endcomment}
\newenvironment{history}{\comment}{\endcomment}

% For commutative diagrams we use Xy-pic
\usepackage[all]{xy}

% We use 2cell for 2-commutative diagrams.
\xyoption{2cell}
\UseAllTwocells

% We use multicol for the list of chapters between chapters
\usepackage{multicol}

% This is generally recommended for better output
\usepackage{lmodern}
\usepackage[T1]{fontenc}

% For cross-file-references

% Package for hypertext links:
\usepackage{hyperref}

% For any local file, say "hello.tex" you want to link to please

% Theorem environments.
%
\theoremstyle{plain}
\newtheorem{theorem}[subsection]{Theorem}
\newtheorem{proposition}[subsection]{Proposition}
\newtheorem{lemma}[subsection]{Lemma}

\theoremstyle{definition}
\newtheorem{definition}[subsection]{Definition}
\newtheorem{example}[subsection]{Example}
\newtheorem{exercise}[subsection]{Exercise}
\newtheorem{situation}[subsection]{Situation}

\theoremstyle{remark}
\newtheorem{remark}[subsection]{Remark}
\newtheorem{remarks}[subsection]{Remarks}

\numberwithin{equation}{subsection}

% Macros
%
\def\lim{\mathop{\mathrm{lim}}\nolimits}
\def\colim{\mathop{\mathrm{colim}}\nolimits}
\def\Spec{\mathop{\mathrm{Spec}}}
\def\Hom{\mathop{\mathrm{Hom}}\nolimits}
\def\Ext{\mathop{\mathrm{Ext}}\nolimits}
\def\SheafHom{\mathop{\mathcal{H}\!\mathit{om}}\nolimits}
\def\SheafExt{\mathop{\mathcal{E}\!\mathit{xt}}\nolimits}
\def\Sch{\mathit{Sch}}
\def\Mor{\mathop{\mathrm{Mor}}\nolimits}
\def\Ob{\mathop{\mathrm{Ob}}\nolimits}
\def\Sh{\mathop{\mathit{Sh}}\nolimits}
\def\NL{\mathop{N\!L}\nolimits}
\def\CH{\mathop{\mathrm{CH}}\nolimits}
\def\proetale{{pro\text{-}\acute{e}tale}}
\def\etale{{\acute{e}tale}}
\def\QCoh{\mathit{QCoh}}
\def\Ker{\mathop{\mathrm{Ker}}}
\def\Im{\mathop{\mathrm{Im}}}
\def\Coker{\mathop{\mathrm{Coker}}}
\def\Coim{\mathop{\mathrm{Coim}}}

% Boxtimes
%
\DeclareMathSymbol{\boxtimes}{\mathbin}{AMSa}{"02}

%
% Macros for moduli stacks/spaces
%
\def\QCohstack{\mathcal{QC}\!\mathit{oh}}
\def\Cohstack{\mathcal{C}\!\mathit{oh}}
\def\Spacesstack{\mathcal{S}\!\mathit{paces}}
\def\Quotfunctor{\mathrm{Quot}}
\def\Hilbfunctor{\mathrm{Hilb}}
\def\Curvesstack{\mathcal{C}\!\mathit{urves}}
\def\Polarizedstack{\mathcal{P}\!\mathit{olarized}}
\def\Complexesstack{\mathcal{C}\!\mathit{omplexes}}
% \Pic is the operator that assigns to X its picard group, usage \Pic(X)
% \Picardstack_{X/B} denotes the Picard stack of X over B
% \Picardfunctor_{X/B} denotes the Picard functor of X over B
\def\Pic{\mathop{\mathrm{Pic}}\nolimits}
\def\Picardstack{\mathcal{P}\!\mathit{ic}}
\def\Picardfunctor{\mathrm{Pic}}
\def\Deformationcategory{\mathcal{D}\!\mathit{ef}}

\setlength{\emergencystretch}{3em}
\begin{document}
\title[Stacks Project, Tag 0D1U]{1270 --- Projective-dimension bounds persist under well-ordered module filtrations over an arbitrary ring}
\maketitle
\noindent Copyright (C) 2005--2025 Johan de Jong. Stacks Project authors. Source: \href{https://stacks.math.columbia.edu/tag/0D1U}{Tag 0D1U}.

\noindent Editorial difficulty note (not author text). Difficulty H1: The proof constructs compatible free modules on the actual elements of each filtration stage, controls the induced kernel filtration at limit ordinals, and combines a transfinite split base case with dimension shifting. This uniform transfinite extension theorem lies beyond ordinary qualification homological algebra..

\noindent Editorial provenance note (not author text). History: excerpt from revision \texttt{a0444\allowbreak 6e57e\allowbreak c1fbc\allowbreak 252a8\allowbreak 71afc\allowbreak ec775\allowbreak 2fb28\allowbreak 07b14}, algebra.tex, lines 26799--26853. Reference presentation was adapted; original-author context and core support blocks were retained or added. The mathematical wording is unchanged. Licensed under GNU FDL 1.2 or later; no invariant sections or cover texts. See ../licenses/COPYING. Context blocks reproduce author definitions and supporting results; they are not additional counted problems.

\begin{lemma}
\label{lemma-exact-sequence-projective-dimension}
Let $R$ be a ring. Let $0 \to M' \to M \to M'' \to 0$ be a short
exact sequence of $R$-modules.
\begin{enumerate}
\item If $M$ has projective dimension $\leq n$ and $M''$
has projective dimension $\leq n + 1$, then $M'$ has projective
dimension $\leq n$.
\item If $M'$ and $M''$ have projective dimension
$\leq n$ then $M$ has projective dimension $\leq n$.
\item If $M'$ has projective dimension $\leq n$ and
$M$ has projective dimension $\leq n + 1$ then
$M''$ has projective dimension $\leq n + 1$.
\end{enumerate}
\end{lemma}

\begin{proof}
Combine the characterization of projective dimension in
Lemma \href{https://stacks.math.columbia.edu/tag/065R}{lemma projective dimension ext (Tag 065R)}
with the long exact sequence of ext groups in
Lemma \href{https://stacks.math.columbia.edu/tag/065P}{lemma reverse long exact seq ext (Tag 065P)}.
\end{proof}


\begin{lemma}
\label{lemma-colimit-projective-dimension}
Let $R$ be a ring. Suppose we have a module $M = \bigcup_{e \in E} M_e$
where the $M_e$ are submodules well-ordered by inclusion. Assume the quotients
$M_e/\bigcup\nolimits_{e' < e} M_{e'}$ have projective dimension $\leq n$.
Then $M$ has projective dimension $\leq n$.
\end{lemma}

\begin{proof}
We will prove this by induction on $n$.

\medskip\noindent
Base case: $n = 0$. Then $P_e = M_e/\bigcup_{e' < e} M_{e'}$ is projective.
Thus we may choose a section $P_e \to M_e$ of the projection $M_e \to P_e$.
We claim that the induced map $\psi : \bigoplus_{e \in E} P_e \to M$ is an
isomorphism. Namely, if $x = \sum x_e \in \bigoplus P_e$ is nonzero,
then we let $e_{max}$ be maximal such that $x_{e_{max}}$ is nonzero
and we conclude that $y = \psi(x) = \psi(\sum x_e)$ is nonzero because
$y \in M_{e_{max}}$ has nonzero image $x_{e_{max}}$ in $P_{e_{max}}$.
On the other hand, let $y \in M$. Then $y \in M_e$ for some $e$.
We show that $y \in \Im(\psi)$ by transfinite induction on $e$.
Let $x_e \in P_e$ be the image of $y$. Then
$y - \psi(x_e) \in \bigcup_{e' < e} M_{e'}$.
By induction hypothesis we conclude that $y - \psi(x_e) \in \Im(\psi)$
hence $y \in \Im(\psi)$. Thus the claim is true and
$\psi$ is an isomorphism. We conclude that $M$ is projective as
a direct sum of projectives, see
Lemma \href{https://stacks.math.columbia.edu/tag/065Q}{lemma direct sum projective (Tag 065Q)}.

\medskip\noindent
If $n > 0$, then for $e \in E$ we denote $F_e$ the free $R$-module
on the set of elements of $M_e$. Then we have a system of
short exact sequences
$$
0 \to K_e \to F_e \to M_e \to 0
$$
over the well-ordered set $E$. Note that the transition maps
$F_{e'} \to F_e$ and $K_{e'} \to K_e$ are injective too.
Set $F = \bigcup F_e$ and $K = \bigcup K_e$. Then
$$
0 \to
K_e/\bigcup\nolimits_{e' < e} K_{e'} \to
F_e/\bigcup\nolimits_{e' < e} F_{e'} \to
M_e/\bigcup\nolimits_{e' < e} M_{e'} \to 0
$$
is a short exact sequence of $R$-modules too and
$F_e/\bigcup_{e' < e} F_{e'}$ is the free $R$-module on the
set of elements in $M_e$ which are not contained in $\bigcup_{e' < e} M_{e'}$.
Hence by
Lemma \ref{lemma-exact-sequence-projective-dimension}
we see that the projective dimension of $K_e/\bigcup_{e' < e} K_{e'}$
is at most $n - 1$. By induction we conclude that $K$ has projective
dimension at most $n - 1$. Whence $M$ has projective dimension at most
$n$ and we win.
\end{proof}
\end{document}
